\documentclass[12pt,a4paper]{article}
\usepackage[margin=2.5cm]{geometry}
\usepackage{amsmath,amssymb}
\usepackage{enumitem}
\setlength{\parindent}{0pt}
\setlength{\parskip}{6pt}

\begin{document}

\begin{center}
\textbf{JAWABAN TUGAS TUTORIAL KE-3}\\
\end{center}

\vspace{2mm}
Nama : RIZIQ LILI ULIL ALBAB \\
NIM \hspace{2.5mm}: 058487469 \\
PRODI : S1 STATISTIKA \\

\vspace{10mm}

%=====================================================
\textbf{1. Tentukan hasil integral berikut.}

\textbf{a.} $\displaystyle \int_{e}^{\infty} \frac{1}{x\ln x}\,dx$

\textit{Jawab:}

Kita akan menggunakan integral tak wajar:
\[
\int_{e}^{\infty} \frac{1}{x\ln x}\,dx = \lim_{b\to\infty} \int_{e}^{b} \frac{1}{x\ln x}\,dx
\]

Misal $u = \ln x$, maka $du = \dfrac{1}{x}\,dx$.
\[
\int \frac{1}{x\ln x}\,dx = \int \frac{1}{u}\,du = \ln|u| + C = \ln|\ln x| + C
\]

Sehingga
\begin{align*}
\int_{e}^{\infty} \frac{1}{x\ln x}\,dx
&= \lim_{b\to\infty} \Big[\ln(\ln x)\Big]_{e}^{b} \\
&= \lim_{b\to\infty} \big(\ln(\ln b) - \ln(\ln e)\big) \\
&= \lim_{b\to\infty} \big(\ln(\ln b) - \ln 1\big) \\
&= \lim_{b\to\infty} \ln(\ln b) - 0 \\
&= \infty
\end{align*}

Integral $\displaystyle \int_{e}^{\infty} \frac{1}{x\ln x}\,dx$ \textbf{divergen}.

\vspace{4mm}
%-----------------------------------------------------
\textbf{b.} $\displaystyle \int_{-\infty}^{\infty} \frac{1}{x^2+16}\,dx$

\textit{Jawab:}

Integral dipecah di $x = 0$:
\[
\int_{-\infty}^{\infty} \frac{1}{x^2+16}\,dx
= \int_{-\infty}^{0} \frac{1}{x^2+16}\,dx + \int_{0}^{\infty} \frac{1}{x^2+16}\,dx
\]

Gunakan rumus $\displaystyle \int \frac{1}{x^2+a^2}\,dx = \frac{1}{a}\arctan\frac{x}{a} + C$ dengan $a = 4$:
\[
\int \frac{1}{x^2+16}\,dx = \frac{1}{4}\arctan\frac{x}{4} + C
\]

Bagian kanan:
\begin{align*}
\int_{0}^{\infty} \frac{1}{x^2+16}\,dx
&= \lim_{b\to\infty} \left[\frac{1}{4}\arctan\frac{x}{4}\right]_{0}^{b} \\
&= \lim_{b\to\infty} \left(\frac{1}{4}\arctan\frac{b}{4} - \frac{1}{4}\arctan 0\right) \\
&= \frac{1}{4}\cdot\frac{\pi}{2} - 0 = \frac{\pi}{8}
\end{align*}

Bagian kiri:
\begin{align*}
\int_{-\infty}^{0} \frac{1}{x^2+16}\,dx
&= \lim_{a\to-\infty} \left[\frac{1}{4}\arctan\frac{x}{4}\right]_{a}^{0} \\
&= \lim_{a\to-\infty} \left(\frac{1}{4}\arctan 0 - \frac{1}{4}\arctan\frac{a}{4}\right) \\
&= 0 - \frac{1}{4}\left(-\frac{\pi}{2}\right) = \frac{\pi}{8}
\end{align*}

Jadi
\[
\int_{-\infty}^{\infty} \frac{1}{x^2+16}\,dx = \frac{\pi}{8} + \frac{\pi}{8} = \frac{\pi}{4}
\]

Integral \textbf{konvergen} dengan nilai $\dfrac{\pi}{4}$.

\vspace{6mm}
%=====================================================
\textbf{2. Tentukan kekonvergenan barisan berikut.}

\textbf{a.} $\displaystyle a_n = \frac{n^3+3n^2+3n}{(n+1)^3}$

\textit{Jawab:}

Karena $(n+1)^3 = n^3+3n^2+3n+1$, sehingga
\[
n^3+3n^2+3n = (n+1)^3 - 1
\]

Maka
\[
a_n = \frac{(n+1)^3 - 1}{(n+1)^3} = 1 - \frac{1}{(n+1)^3}
\]

Sehingga
\[
\lim_{n\to\infty} a_n = \lim_{n\to\infty}\left(1 - \frac{1}{(n+1)^3}\right) = 1 - 0 = 1
\]

Barisan $\{a_n\}$ \textbf{konvergen} ke 1.

\vspace{4mm}
%-----------------------------------------------------
\textbf{b.} $\displaystyle b_n = \left(1+\frac{2}{n}\right)^{n/2}$

\textit{Jawab:}

$1^\infty$ (bentuk tak tentu). Misalkan
\[
y = \left(1+\frac{2}{x}\right)^{x/2}, \quad x > 0
\]

Ambil logaritma natural kedua ruas:
\[
\ln y = \frac{x}{2}\ln\left(1+\frac{2}{x}\right) = \frac{\ln\left(1+\frac{2}{x}\right)}{\frac{2}{x}}
\]

$x\to\infty$ bentuknya $\dfrac{0}{0}$. Misalkan $t = \dfrac{2}{x}$, maka $x\to\infty \Rightarrow t\to 0^{+}$:
\[
\lim_{x\to\infty} \ln y = \lim_{t\to 0^{+}} \frac{\ln(1+t)}{t}
\]

Pakai aturan L'H\^opital:
\[
\lim_{t\to 0^{+}} \frac{\ln(1+t)}{t} = \lim_{t\to 0^{+}} \frac{\frac{1}{1+t}}{1} = \frac{1}{1+0} = 1
\]

Karena $\ln y \to 1$, maka $y \to e^{1} = e$. Jadi
\[
\lim_{n\to\infty} b_n = \lim_{n\to\infty}\left(1+\frac{2}{n}\right)^{n/2} = e
\]

Barisan $\{b_n\}$ \textbf{konvergen} ke $e$.

\vspace{6mm}
%=====================================================
\textbf{3. Tentukan kekonvergenan deret berikut.}

\textbf{a.} $\displaystyle \sum_{k=1}^{\infty} \frac{4^{k+1}}{7^{k-1}}$

\textit{Jawab:}

Ubah suku ke-$k$ ke bentuk deret geometri:
\[
\frac{4^{k+1}}{7^{k-1}} = \frac{4^{2}\cdot 4^{k-1}}{7^{k-1}} = 16\left(\frac{4}{7}\right)^{k-1}
\]

Sehingga
\[
\sum_{k=1}^{\infty} \frac{4^{k+1}}{7^{k-1}} = \sum_{k=1}^{\infty} 16\left(\frac{4}{7}\right)^{k-1}
\]

Deret geometri dengan suku pertama $a = 16$ dan rasio $r = \dfrac{4}{7}$.

Karena $|r| = \dfrac{4}{7} < 1$, deret konvergen dengan jumlah
\[
S = \frac{a}{1-r} = \frac{16}{1-\frac{4}{7}} = \frac{16}{\frac{3}{7}} = \frac{112}{3}
\]

Deret \textbf{konvergen} dengan jumlah $\dfrac{112}{3}$.

\vspace{4mm}
%-----------------------------------------------------
\textbf{b.} $\displaystyle \sum_{k=2}^{\infty} \left(\frac{3}{(k-1)^2} - \frac{3}{k^2}\right)$

\textit{Jawab:}

Deret ini adalah deret kolaps (\textit{telescoping}). Jumlah parsial ke-$n$:
\begin{align*}
S_n &= \sum_{k=2}^{n} \left(\frac{3}{(k-1)^2} - \frac{3}{k^2}\right) \\
&= \left(\frac{3}{1^2} - \frac{3}{2^2}\right) + \left(\frac{3}{2^2} - \frac{3}{3^2}\right) + \left(\frac{3}{3^2} - \frac{3}{4^2}\right) + \cdots + \left(\frac{3}{(n-1)^2} - \frac{3}{n^2}\right) \\
&= \frac{3}{1^2} - \frac{3}{n^2} \\
&= 3 - \frac{3}{n^2}
\end{align*}

Maka
\[
\sum_{k=2}^{\infty} \left(\frac{3}{(k-1)^2} - \frac{3}{k^2}\right) = \lim_{n\to\infty} S_n = \lim_{n\to\infty}\left(3 - \frac{3}{n^2}\right) = 3 - 0 = 3
\]

Deret \textbf{konvergen} dengan jumlah 3.

\end{document}
